\section{方差、矩与尺度}
\label{sec:05-Probability/04-Expectation-and-Moments/02}

返回到上一节结尾的那对投资品。两个产品承诺的年化期望收益率都是 10\%。产品 A 很简单——每年稳定回报 10\%。产品 B 呢？一半年份涨 100\%，一半年份跌 80\%。算一下期望：\(0.5\times1.0+0.5\times(-0.8)=0.10\)，确实也是 10\%。可你不会说二者没区别——产品 B 让你有一半概率接近血本无归。

期望只看中心，不碰风险。这一节要找一个描述 ``分散程度'' 的量。

\subsection{直接平均偏差为什么不行}

最自然的念头：把每次偏离均值的距离直接平均：

\[
\E[X-\mu],\qquad \mu=\E[X].
\]

试试看。公平骰子：\(\mu=3.5\)。偏离分别是 \(-2.5,-1.5,-0.5,0.5,1.5,2.5\)，平均值恰好是 0。不是巧合——正负偏差在取期望时必然对消，因为均值本就定义在分布重心。这个量永远为零，不能区分任何分布。

那就取绝对值：\(\E[|X-\mu|]\)。它确实能量化分散，且不放大极端值。问题在于绝对值函数在零点不可导，解析操作困难。对它求和时无法把变量拆开（两个变量和的平均绝对偏差不等于各自平均绝对偏差之和）。

平方偏离给出了更好的选择：

\[
\Var(X)=\E[(X-\mu)^2].
\]

平方消灭了正负号，可导，而且对大偏离给了不成比例的惩罚——直觉上正好匹配我们对 ``风险'' 的感受：偏离两倍比偏离一倍糟了不止两倍。

\subsection{展开式提供计算捷径}

把平方展开：

\[
\begin{aligned}
\Var(X)&=\E[X^2-2\mu X+\mu^2]\\
&=\E[X^2]-2\mu\,\E[X]+\mu^2\\
&=\E[X^2]-2\mu^2+\mu^2\\
&=\E[X^2]-(\E[X])^2.
\end{aligned}
\]

这是纯代数变形，不需要新假设。多数情况直接用这个公式算方差比回到定义更快。

用公平骰子走一遍数值。我们已经知道（上一节用 LOTUS）：

\[
\E[X]=\frac72=3.5,\qquad
\E[X^2]=\frac{91}{6}\approx15.167.
\]

代入：

\[
\Var(X)=\frac{91}{6}-\Bigl(\frac72\Bigr)^2
=\frac{91}{6}-\frac{49}{4}
=\frac{182-147}{12}
=\frac{35}{12}\approx2.917.
\]

数字自己跑完。你会注意到 \(\Var(X)\ge0\)——这来自平方的非负性，\(E[X^2]\ge (\E[X])^2\) 正是 Jensen 不等式在 \(x\mapsto x^2\) 上的特例。

方差什么时候等于零？只有当 \(X=\mu\) 几乎必然成立时。如果 \(X\) 有任何概率落在离 \(\mu\) 大于 \(\varepsilon>0\) 的地方，它对平方平均的贡献至少是 \(\varepsilon^2 P(|X-\mu|>\varepsilon)>0\)。所以零方差等于没有随机性——分布退化成一个点。

\subsection{标准差：把单位从平方拉回来}

方差的单位是原变量单位的平方。身高用厘米计时，方差单位是 \(\text{cm}^2\)，没法跟原始数据直接比较。标准差解决了这个问题：

\[
\sigma_X=\sqrt{\Var(X)}.
\]

标准差与 \(X\) 同单位。骰子的 \(\sigma=\sqrt{35/12}\approx1.708\)——单次掷骰结果与 3.5 的典型偏离大约在 1.7 点左右。

线性变换时量纲变化很清楚：

\[
\E[aX+b]=a\E[X]+b,\qquad
\Var(aX+b)=a^2\Var(X).
\]

加上常数 \(b\) 只是整条分布左右平移，散布程度不变——方差和标准差都对平移免疫。乘上常数 \(a\) 则把所有偏离放大 \(|a|\) 倍，方差因此放大 \(a^2\) 倍。

这引出一个实用的技巧——标准化：

\[
Z=\frac{X-\mu}{\sigma},\qquad(\sigma>0).
\]

则 \(\E[Z]=0,\;\Var(Z)=1\)。不同物理单位的变量现在可以放在同一尺度上比较： ``比均值高出两个标准差'' 无论是来自于身高、收入还是考试成绩，含义是相通的。

\subsection{高阶矩：不只关心散布}

方差是二阶中心矩。一般地，\(k\) 阶原点矩是 \(\E[X^k]\)，\(k\) 阶中心矩是 \(\E[(X-\mu)^k]\)。

一阶原点矩是均值，二阶中心矩是方差。三阶中心矩标准化后成为偏度（skewness），描述分布的不对称方向：正偏意味着右尾比左尾更长，均值在众数右侧。四阶标准化中心矩关联峰度（kurtosis），描述尾部的厚薄——大峰度意味着分布有更 ``重'' 的尾巴，极端值比正态分布的预期更频繁。

但这里需要一句冷静的话：知道再多矩也不能代替知道分布本身。

\begin{itemize}
\item 高阶矩可能根本不存在。上一节那个 \(P(X=2^k)=2^{-k}\) 的例子连一阶矩都发散，还谈什么三阶。
\item 有限个矩从不唯一确定分布。两个不同的分布可以有相同的前 \(n\) 阶矩。
\item 即使所有阶矩都存在，也不保证矩唯一确定分布（需要满足 Carleman 条件之类额外假设）。
\end{itemize}

矩是分布的素描，不是照片。

\subsection{均方误差：用预测的视角理解均值与方差}

用一个常数 \(a\) 预测 \(X\)，自然用平方误差 \((X-a)^2\) 衡量预测质量，期望误差为

\[
\E[(X-a)^2].
\]

把这拆开，加减 \(\mu=\E[X]\)：

\[
\begin{aligned}
\E[(X-a)^2]&=\E[(X-\mu+\mu-a)^2]\\
&=\E[(X-\mu)^2]+2(\mu-a)\underbrace{\E[X-\mu]}_{=0}+(\mu-a)^2\\
&=\Var(X)+(\mu-a)^2.
\end{aligned}
\]

交叉项消掉是因为偏离均值的期望为零。两项各自有明确的含义：

\begin{itemize}
\item \(\Var(X)\)：即使用最优常数，仍然剩下来的随机波动——这部分你无法通过选择 \(a\) 消除。
\item \((\mu-a)^2\)：预测常数离均值有多远——离得越远，损失越大。
\end{itemize}

因此，在平方损失下，最佳常数预测就是 \(a=\E[X]\)，此时达到的最小均方误差等于 \(\Var(X)\)。这给出了均值和方差一个不同于 ``长期平均'' 的诠释：均值是使预测平方误差最小的那个数，方差是用了最好预测之后仍然无法避免的残余误差。下一节再遇到 ``条件期望是最优预测'' 的时候，这套逻辑会往条件方向再走一步。

\subsection{方差对异常值太过忠诚}

平方放大了距离。一个偏离均值 10 倍标准差的点对平方平均的贡献是 100 个单位。如果数据由少量极端值和大量集中值混合构成，方差会主要反映那些极端值的特征，而不是 ``典型'' 数据的行为。

更严重的问题是：有限均值不担保有限方差。取

\[
P(X=k)=\frac{c}{k^3},\qquad k=1,2,\ldots
\]

其中 \(c\) 是归一化常数使概率和为 1。则

\[
\E[X]=c\sum_{k=1}^{\infty}\frac1{k^2}=c\cdot\frac{\pi^2}{6}<\infty,
\]
但
\[
\E[X^2]=c\sum_{k=1}^{\infty}\frac1{k}\to\infty.
\]

分布有有限均值却没有有限方差。实际中遇到这类重尾时，样本方差可能不收敛到一个稳定值——增加数据不能让它安分下来。这时需要考虑评估波动性的其他度量（分位数间距、平均绝对偏差等），或者换用适合重尾的损失函数。

方差是一个好用的默认工具，但它假设你的损失函数确实接近平方、你的尾部确实不太重。如果这两个前提不成立，需要去工具箱里翻别的器械。
